In this chapter I test whether the models fine tuned on the survey data respond differently to various moral, ethical and value dilemmas. The goal here is to firstly answer \rqref{align} by comparing base models to fine-tuned models and then secondly use the model responses in \ref{chap:survey} to answer \rqref{prefer}. I do this by presenting models with self contained dilemmas across a wide range of NZ work contexts (namely public sector jobs). The model is asked to provide a response which describes what their decision would be. In this way we can compare the responses for differences between base models and fine-tuned models. Initially the methodology was going to test the model using agentic dilemmas, however this proved too difficult to get clean differences of value judgements. Therefore a soft pivot was conducted to use single response dilemmas instead. More information about the agentic dilemmas can be found in \ref{app:simulation}

The structure of the chapter is as follows:

\section{The difference between base and fine-tuned models}

This is the results section which goes first.

\subsection{Situational awareness of dilemmas}

Models are shown to know that they are in a simulation \cite{laineMeMyselfAI2024,phuongEvaluatingFrontierModels2025,needhamLargeLanguageModels2025}. Smaller models areless self aware \cite{chaudharyEvaluationAwarenessScales2025}

\section{Creating a small Dilemma dataset by hand}

The dilemmas were constructed by hand with three objectives in in mind. These objectives exist to make the project tractable and to ensure that the downstream survey is of high quality. The objectives are as follows:

\begin{enumerate}
    \item Each dilemma should be self contained and simple enough that any New Zealander could read and understand it without needing to know the context of the project.
    \item The dilemma should reveal a clear decision that must be made.
    \item Each of the decisions that could be made must be defensible under some commonsense moral or ethical framework.
\end{enumerate}

Notably the value tensions that are tested is not a comprehensive list of all value tensisons and instead is a best-effort coverage of various value tensions.


\subsection{Writing of the dilemmas}
It was an interactive design process in which initially 10 dilemmas were constructed. 
Then these ten dilemmas were given two variations each which target a similar trade off.

Iterativing with teesting with a ensemble of frontier models to test fo balance, divergence and stabilty

\section{Dilemmas as single response agents}
The goal of the Dilemmas was to test the models responses in an out of distriubtion setting.

This mean that maintaining some realism is required to invoke the model to respond 'as if' it were in deployment.

The way to do this was the framing of the system prompt in which each model was given a systme pormt htat both told them they were the first step in an agentic process as well as giving them a large amount of relevant documents to reference (e.g Code of Conducts). These reference documents were real public documents where relevant.

\section{How the dilemmas were graded}

They were graded using LLM as Judge. Wuith the judge being asked two things. To decide which of a defined set of options were taken and then to provide a deciion of which value it prioritised.

\ctable{dilemmas-answer-stats}{The answer and stability satistics for the dilemmas. The mean answer is the average of the model responses across the latest run set, and the stabilty is the entropy of the model responses. The bolded values indicate the highest stability for each dilemma.}

\ctable{dilemmas-answer-stats-frontier}{The same answer and stability statistics for the frontier comparison models the dilemmas were also run on (openrouter; the presenter and judge were unchanged).}
